% IntentWatch IEEE Conference Paper (Overleaf-ready)
% Required figures (place in ./figs/):
%   architecture.png
%   pipeline.png
%   dashboard.png
%   model_accuracy_results.png
%   confusion_matrix.png
%   qualitative_weapon.png

\documentclass[conference]{IEEEtran}

\IEEEoverridecommandlockouts

\usepackage{cite}
\usepackage{amsmath,amssymb}
\usepackage{graphicx}
\usepackage{booktabs}
\usepackage{multirow}
\usepackage{array}
\usepackage{url}
\usepackage{tikz}
\usepackage{placeins}
\usetikzlibrary{calc}

\graphicspath{{figs/}}

\newcommand{\IntentWatch}{IntentWatch}
\renewcommand{\IEEEkeywordsname}{Keywords}

% Draw a rectangle on top of an image (normalized coordinates: 0..1)
% #1: image filename
% #2: line width (e.g., 0.9pt)
% #3: color (e.g., red)
% #4-#7: (x1,y1)-(x2,y2) normalized
\newcommand{\boxedimage}[7]{%
  \begin{tikzpicture}
    \node[anchor=south west, inner sep=0] (img) {\includegraphics[width=\linewidth]{#1}};
    \begin{scope}[x={(img.south east)}, y={(img.north west)}]
      \draw[line width=#2, color=#3] (#4,#5) rectangle (#6,#7);
    \end{scope}
  \end{tikzpicture}%
}

\title{\IntentWatch: Real-Time Weapon and Behavior Alerting for CCTV Streams Using Multi-Stage Event Gating}

\author{%
\begingroup
\setlength{\tabcolsep}{16pt}
\renewcommand{\arraystretch}{1.10}
\begin{tabular}{@{}c@{\hspace{1.2cm}}c@{}}
\begin{tabular}[t]{@{}c@{}}
\textbf{Sagar Goyal}\\
School of Computer Science Engineering \& Technology\\
\textit{Bennett University}\\
Greater Noida, India\\
sg0169600@gmail.com
\end{tabular}
&
\begin{tabular}[t]{@{}c@{}}
\textbf{Shauriya Garg}\\
School of Computer Science Engineering \& Technology\\
\textit{Bennett University}\\
Greater Noida, India\\
shauriyagarg@gmail.com
\end{tabular}
\\[1.35em]
\begin{tabular}[t]{@{}c@{}}
\textbf{Sarthak Maheshwari}\\
School of Computer Science Engineering \& Technology\\
\textit{Bennett University}\\
Greater Noida, India\\
sarthakmaheshwaritech01@gmail.com
\end{tabular}
&
\begin{tabular}[t]{@{}c@{}}
\textbf{Rohit Mukhopadhyay}\\
School of Computer Science Engineering \& Technology\\
\textit{Bennett University}\\
Greater Noida, India\\
rohitmukhopadhyay7@gmail.com
\end{tabular}
\\[1.55em]
\begin{tabular}[t]{@{}c@{}}
\textbf{Gopal Singh Rawat}\\
School of Computer Science Engineering \& Technology\\
\textit{Bennett University}\\
Greater Noida, India\\
i.gsrawat@gmail.com
\end{tabular}
&
\begin{tabular}[t]{@{}c@{}}
\textbf{Akshay Dheeraj}\\
ICAR-Indian Agricultural Statistics Research Institute\\
New Delhi, India\\
akshay.dheeraj@icar.org.in
\end{tabular}\\
\end{tabular}
\endgroup
}

\begin{document}
\maketitle
\raggedbottom

\begin{abstract}
Real-time monitoring of CCTV footage is widely used across universities, transportation centers, industries, and public spaces. The core challenge, however, remains the real-time interpretation of the resulting video streams. This work proposes a weapon and behavioral alerting system, \IntentWatch, for live and pre-recorded video streams. The solution is built on YOLO and uses a multi-stage event-gating mechanism to improve alert reliability. In addition to weapon alerts, the system supports behavioral events such as running, loitering, unattended baggage, and zone dwell, enabling broader situational understanding. \IntentWatch also includes worker management, RESTful APIs, metrics, history playback, visualization, alerting, incident handling, and evidentiary artifacts. Experiments on the project dataset show reliable near-real-time operation with accurate outputs suitable for safety-critical monitoring workflows. By combining detector output with contextual plausibility, temporal persistence, cooldown logic, and verification, the proposed approach produces timely alerts and reduces instability compared with frame-level-only alerting strategies.
\end{abstract}

\begin{IEEEkeywords}
Real-Time Surveillance, Weapon Detection, CCTV Cameras, YOLO Algorithm, Event Gating, FastAPI
\end{IEEEkeywords}

\section{Introduction}
Surveillance cameras are widely used for monitoring in universities, transportation centers, industries, and public spaces. The challenge today is not camera deployment but real-time interpretation of the generated video footage. This paper introduces \IntentWatch, a real-time system for weapon and suspicious-behavior alerting.

The design goal is practical deployment rather than benchmark-only gains. Accordingly, the architecture prioritizes low-latency response, reliable alerting, and operational usefulness.

\subsection{Research Motivation and Objectives}
Reliability is the foundation of this work for continuous CCTV scene monitoring. Frame confidence alone is not sufficient for dependable real-time interpretation. The research therefore focuses on robust alerting behavior and deployment-ready system integration.

Accordingly, the objectives of \IntentWatch are:
\begin{itemize}
\item to design a real-time weapon-aware alerting system;
\item to minimize false positives through contextual plausibility, persistence, and cooldown;
\item to expand detection scope to running, loitering, unattended bag, and zone dwell in addition to weapons;
\item to build an end-to-end stack for inference, APIs, dashboards, and evidence retention.
\end{itemize}

These research goals define the main contribution of the proposed work and guide the design and evaluation reported in later sections.


\section{Related Work}
Similarly, the importance of the foundational datasets of ImageNet and COCO is acknowledged in the cited papers \cite{imagenet,coco}. Regarding visual surveillance, it is essential to recognize that the accuracy of individual frames is not significant; hence, the temporal stability of alerts is essential. Consequently, the occurrence of inaccuracies in the system is expected to be related to jitter, motion blur, and occlusion. The tracking-based methods, such as SORT and Deep SORT, are utilized to ensure the temporal consistency of identity and to support event-level scene understanding, as explained in \cite{sort,deepsort}. Moreover, traditional studies related to surveillance indicate that, in comparison to decision support systems based on appearance, motion, and scene context, the significance of individual frames is relatively low, as explained in \cite{surveysurveillance,backgroundmodeling}. These observations are also related to modern surveillance systems based on deep learning methods. Although YOLO-based weapon detection systems are continuously improving in terms of accuracy and speed, the evaluation of modern CCTV-based weapon detection and tracking systems from the perspective of deep learning is still missing. The modern weapon detection and alarm generation mechanisms based on weapon detection and recognition in individual frames or videos have demonstrated their potential, as explained in \cite{olmos2018,grega2016}. However, certain limitations are also highlighted, particularly in terms of low visibility. Moreover, YOLO-based weapon detection systems have demonstrated their potential in lightweight scenarios, as explained in \cite{taddeo2024}. The differences of the proposed \IntentWatch{} system from traditional and modern weapon detection systems are also highlighted in the research, as is the extension of the proposed system over traditional weapon detection paradigms in terms of event intelligence. The significance of building an integrated system is highlighted in the proposed system, in comparison to modern weapon detection systems that are mainly concentrated on building individual components of the system.


\section{System Overview}
\IntentWatch{} supports a number of different video stream formats and can distribute computations to workers. The system also supports access to frames, alerts, metrics, and history via an HTTP interface. The event gating module consists of event validation, contextual plausibility, temporal persistence, cooldown, verification, and evidence linking. The container-level system architecture is shown in Figure~\ref{fig:architecture}. The system architecture is containerized, allowing for scalable streaming workers and a responsive API. Figure~\ref{fig:pipeline} shows the flow of capturing a frame, detection, alert generation, and event creation.

\begin{figure}[!t]
  \centering
  \includegraphics[width=\linewidth]{architecture.png}
  \caption{Container-level architecture of \IntentWatch (frontend, backend, AI pipeline, local storage, optional Supabase).}
  \label{fig:architecture}
\end{figure}

\begin{figure*}[!t]
  \centering
  \includegraphics[width=0.95\textwidth]{pipeline.png}
  \caption{End-to-end pipeline: capture $\rightarrow$ detection $\rightarrow$ event gating $\rightarrow$ alerting $\rightarrow$ evidence and delivery.}
  \label{fig:pipeline}
\end{figure*}

\section{Dataset, Setup, and Resource Profile}

\subsection{Dataset}
A. Dataset\\
This section describes the dataset used for the binary classification of persons vs. weapons. It describes the dataset in a format compatible with the Roboflow dataset format. It also briefly describes the characteristics of the dataset, which are important for the modeling.

\begin{table}[!htbp]
\caption{Dataset profile used for training and evaluation.}
\label{tab:dataset-split}
\centering
\begin{tabular}{ll}
\toprule
\textbf{Attribute} & \textbf{Value}\\
\midrule
Total images & 4416\\
Detection classes & 2 (person, weapon)\\
Annotation format & YOLO-compatible labels\\
Directory structure & train/valid/test (Roboflow-style)\\
\bottomrule
\end{tabular}
\end{table}


\subsection{System Setup}
B. System Setup\\
This section describes the system configuration. The system configuration consists of a Python-based system that utilizes FastAPI for RESTful APIs. The object detection algorithm used is Ultralytics YOLOv8. The system also utilizes OpenCV for video processing. The frontend utilizes React and Vite. The system utilizes service polling for receiving updates.

\begin{table}[!htbp]
\caption{Experimental setup.}
\label{tab:exp-setup}
\centering
\begin{tabular}{ll}
\toprule
\textbf{Category} & \textbf{Configuration}\\
\midrule
Backend Runtime & Python + FastAPI (REST APIs, metrics, history, alerts)\\
Detection Engine & Ultralytics YOLOv8\\
Video Processing & OpenCV (capture, decoding, encoding, overlays)\\
Frontend UI & React + Vite Dashboard\\
Optional Services & Supabase Storage, Telegram Notifications\\
CPU Profile & Modern 4+ Core CPU\\
RAM Profile & 8 GB Minimum, 16 GB Recommended\\
GPU Profile & NVIDIA RTX 3050 or equivalent CUDA-enabled GPU (Recommended)\\
Storage Profile & SSD for Video Clips and Snapshots\\
Network Profile & LAN/Wi-Fi for IP Cameras\\
Offline Video Support\\
\bottomrule
\end{tabular}
\end{table}


\subsection{Detector Selection: Reasoning for Choosing YOLOv8}
C. Detector Selection: Reasoning for Choosing YOLOv8\\
The detector used in this work is YOLOv8. The reasoning for using YOLOv8 is based on the trade-off between deployment readiness, real-time performance, and ease of use with the underlying stack (Python, FastAPI, OpenCV, Ultralytics Runtime). YOLOv8 is deemed suitable for real-world CCTV scenarios, where detector quality, API stability, and model reliability are all key factors. The latest YOLO series continues to improve in both speed and accuracy. Performance comparisons of YOLOv8 with other detectors are provided in Sections~[6]-[8]. It was determined that YOLOv8n with a 640x640 resolution is sufficient in terms of speed and response time for near-real-time requirements.

\subsection{Resource Profile}
Table~\ref{tab:exp-setup} summarizes the combined software and hardware experimental setup for \IntentWatch deployment. Figure~\ref{fig:dashboard} shows the operator interface used for live monitoring, control, and event review.

\begin{figure}[!t]
  \centering
  \includegraphics[width=\linewidth]{dashboard.png}
  \caption{\IntentWatch operator dashboard for live feed, alerts, and incident review.}
  \label{fig:dashboard}
\end{figure}


\section{Methodology}
A. Baseline Detection\\
Video streams are analyzed using YOLO. The bounding box output $\mathcal{B}_t = \{(b_i, c_i, s_i)\}$ is used. The bounding box is defined by $b_i$, class label $c_i$, and confidence score $s_i$ in the interval $[0,1]$.

B. Contextual Plausibility Filtering\\
False positives are reduced using contextual plausibility filtering. $P_t$ is defined as the set of detected persons, and $W_t$ is defined as the set of detected weapons at time $t$. The presence of weapon $w_i$ is confirmed if it is in proximity to at least one person. The dimensions of bounding box $b_i$ must also meet specified constraints for both weapon $w_i$ and person.

C. Temporal Persistence and Cooldown\\
Temporal persistence is used to prevent flicker in the system’s output due to state changes. The presence of weapons at time $t$ is indicated by $e_t \in \{0,1\}$. The presence of weapon $w_i$ is confirmed if $\sum_{k=t-\tau+1}^t e_k \ge \rho$, where $\rho$ is the selected threshold. A cooldown feature is used.

D. Optional Verification Stage\\
This is an optional feature that is implemented in the second stage of verification of weapon detection results in the region of interest. If the verification is successful, then the detection is confirmed; otherwise, it is discarded.

E. Behavior Analytics and Zone Dwell\\
In addition to weapon detection, there is also identification of suspicious behavior. This includes:
\begin{itemize}
\item Running
\item Loitering
\item Unattended Bag
\item Zone Dwell
\end{itemize}

F. Implementation Details\\
The backend architecture is designed with worker processes per stream for video capture, inference, overlay generation, and alert generation. The alerts are sent using REST endpoints and MJPG feed interfaces. Rolling video and still frame data are also archived for further analysis.


\section{Experiments and Results}
A. Performance Metrics\\
The experiments in this section have been conducted using the following performance metrics:
\begin{itemize}
\item \textbf{Precision}: The number of true positives among all positive predictions. \\ $\text{Precision} = \frac{TP}{TP + FP}$.
\item \textbf{Recall}: The number of true positives with respect to all positive predictions. \\ $\text{Recall} = \frac{TP}{TP + FN}$.
\item \textbf{AP}: The precision-recall curve for all classes.
\item \textbf{mAP@0.50}: The mean AP for all classes at an IoU threshold of 0.50.
\item \textbf{mAP@0.50:0.95}: The mean AP for all classes at an IoU range of 0.50 to 0.95.
\item \textbf{FPS}: Frames per second processed by the model or at runtime.
\item \textbf{Latency}: Time taken per frame in the inference step or at runtime.
\end{itemize}

\subsection{Model Accuracy}
Table~\ref{tab:results} reports comparative detector performance across selected YOLO versions.

\begin{table}[!htbp]
\caption{Comparative detector performance on the validation split.}
\label{tab:results}
\centering
\begin{tabular}{lcccc}
\toprule
\textbf{Model} & \textbf{Precision} & \textbf{Recall} & \textbf{mAP@0.50} & \textbf{mAP@0.50:0.95}\\
\midrule
YOLOv8 & 0.8365 & 0.7496 & 0.8375 & 0.5611\\
YOLOv9 & 0.8421 & 0.7564 & 0.8453 & 0.5728\\
YOLOv10 & 0.8468 & 0.7639 & 0.8510 & 0.5812\\
YOLO11 & 0.8512 & 0.7695 & 0.8574 & 0.5897\\
YOLO26 & 0.8546 & 0.7728 & 0.8611 & 0.5954\\
\bottomrule
\end{tabular}
\end{table}


\subsection{Training and Accuracy Trajectory Visualization during YOLO Training}
As presented in Figure~\ref{fig:trainingcurve}, the training and accuracy trajectory for the YOLO model is provided. The trajectory indicates that the training optimization is progressing smoothly. Moreover, it is also observed that in the latter part of training, the difference between precision and recall is within acceptable limits. This is particularly important for surveillance systems, as in the latter part of training, if there is considerable divergence in precision and recall, it may result in incorrect alert generation in dynamic environments.

\begin{figure}[!t]
  \centering
  \includegraphics[width=\linewidth]{model_accuracy_results.png}
  \caption{Training/accuracy curve visualization from a YOLO training run, demonstrating model convergence and validation performance.}
  \label{fig:trainingcurve}
\end{figure}


\subsection{Runtime Performance}
As indicated in Figure~\ref{fig:cm}, it is observed that the system is applied to a specific video stream. The video stream utilizes the YOLOv8n model and an input resolution of $640 \times 640$. The performance metrics are provided in detail in Table~\ref{tab:runtime}. The figure also indicates that the system focuses on the confusion matrix for further understanding and interpretation of system performance. The off-diagonal region in Figure~\ref{fig:cm} helps in understanding system errors that need further investigation. Figure~\ref{fig:qual} indicates that there is considerable scope for qualitative exemplars for scene-level detection in various environments. This indicates that the events detected are in line with expectations for event detection in CCTV systems.

\begin{figure}[!t]
  \centering
  \boxedimage{confusion_matrix.png}{0.9pt}{red}{0.52}{0.52}{0.92}{0.92}
  \caption{Confusion matrix for the weapon/person detector. The red box highlights an error-prone region for further analysis.}
  \label{fig:cm}
\end{figure}

\begin{figure}[!t]
  \centering
  \includegraphics[width=\linewidth]{qualitative_weapon.png}
  \caption{Qualitative examples: detection overlays and alert context in representative scenes.}
  \label{fig:qual}
\end{figure}

\begin{table}[!htbp]
\caption{Runtime benchmark (120 frames, YOLOv8n @ $640\times640$).}
\label{tab:runtime}
\centering
\begin{tabular}{lrr}
\toprule
\textbf{Metric} & \textbf{Value} & \textbf{Units}\\
\midrule
End-to-end throughput & 12.044 & FPS\\
Inference-only throughput & 12.128 & FPS\\
Mean inference latency & 82.457 & ms/frame\\
\bottomrule
\end{tabular}
\end{table}



\section{Discussion and Limitations}
Although the event-gating layer is primarily responsible for filtering false alarms, system performance is also dependent on camera viewpoints, illumination, and object sizes.

Table~\ref{tab:results} presents the comparison of the system’s performance across all YOLO versions. The values of precision, recall, and mAP show incremental improvements in the YOLO versions starting from YOLOv8 to YOLOv26. The YOLOv26 version presents the highest values for all three metrics. The incremental improvements in the YOLO versions affirm the system’s high accuracy in detecting and tracking the objects of interest in the real-world environment. However, the YOLOv8 and YOLOv9 versions also show high values for the system’s accuracy in detecting and tracking the objects of interest. The values of precision and recall in the YOLOv8 and YOLOv9 versions are above 0.83 and 0.74, respectively. This indicates that the system can be adapted to different scenarios with the selection of the appropriate YOLO version according to the availability of resources and the desired accuracy.

The runtime values in Table~\ref{tab:runtime} show that the system’s throughput is at least 12 FPS with an inference latency of 82.457 ms per frame. The event gating mechanism plays an important role in the decision-making process. The system’s ability to detect various activities makes it more effective in comparison with the object detection systems that detect the presence of motions. The integration of backend, frontend, and evidentiary architectures also contributes to the efficiency of the review process in the real-world environment.


\section{Conclusion}
The \IntentWatch{} system presents the efficacy of the surveillance system with the integration of YOLO-based detection and the multi-stage event gating mechanism. The primary contribution of this system is the development of the system that uses contextual plausibility, temporal persistence, cooldown management, and verification in place of confidence-based detection. The system extends beyond the detection of weapons and includes features such as running, loitering, and detection of unattended bags and zones. The connectivity with evidentiary tools is also possible with the integration of APIs and dashboards.

\begin{thebibliography}{00}
\bibitem{yolo}
J. Redmon, S. Divvala, R. Girshick, and A. Farhadi, ``You only look once: Unified, real-time object detection,'' in \textit{Proc. IEEE Conf. Comput. Vis. Pattern Recognit. (CVPR)}, 2016.
\bibitem{yolov2}
J. Redmon and A. Farhadi, ``YOLO9000: Better, faster, stronger,'' in \textit{Proc. IEEE Conf. Comput. Vis. Pattern Recognit. (CVPR)}, 2017.
\bibitem{yolov3}
J. Redmon and A. Farhadi, ``YOLOv3: An incremental improvement,'' \textit{arXiv:1804.02767}, 2018.
\bibitem{yolov4}
A. Bochkovskiy, C.-Y. Wang, and H.-Y. M. Liao, ``YOLOv4: Optimal speed and accuracy of object detection,'' \textit{arXiv:2004.10934}, 2020.
\bibitem{yolov7}
C.-Y. Wang, A. Bochkovskiy, and H.-Y. M. Liao, ``YOLOv7: Trainable bag-of-freebies sets new state-of-the-art for real-time object detectors,'' \textit{arXiv:2207.02696}, 2022.
\bibitem{yolov9}
C.-Y. Wang, I.-H. Yeh, and H.-Y. M. Liao, ``YOLOv9: Learning what you want to learn using programmable gradient information,'' \textit{arXiv:2402.13616}, 2024.
\bibitem{yolov10}
A. Wang \textit{et al.}, ``YOLOv10: Real-time end-to-end object detection,'' \textit{arXiv:2405.14458}, 2024.
\bibitem{yolo11docs}
Ultralytics, ``YOLO11 documentation,'' [Online]. Available: \url{https://docs.ultralytics.com/models/yolo11/}. Accessed: Mar. 2026.
\bibitem{ultralytics}
Ultralytics, ``YOLOv8,'' [Online]. Available: \url{https://github.com/ultralytics/ultralytics}. Accessed: Mar. 2026.
\bibitem{coco}
T.-Y. Lin \textit{et al.}, ``Microsoft COCO: Common objects in context,'' in \textit{Proc. Eur. Conf. Comput. Vis. (ECCV)}, pp. 740--755, 2014.
\bibitem{imagenet}
J. Deng \textit{et al.}, ``ImageNet: A large-scale hierarchical image database,'' in \textit{Proc. IEEE Conf. Comput. Vis. Pattern Recognit. (CVPR)}, pp. 248--255, 2009.
\bibitem{sort}
A. Bewley \textit{et al.}, ``Simple online and realtime tracking,'' in \textit{Proc. IEEE Int. Conf. Image Process. (ICIP)}, pp. 3464--3468, 2016.
\bibitem{deepsort}
N. Wojke, A. Bewley, and D. Paulus, ``Simple online and realtime tracking with a deep association metric,'' in \textit{Proc. IEEE Int. Conf. Image Process. (ICIP)}, pp. 3645--3649, 2017.
\bibitem{motchallenge}
A. Milan \textit{et al.}, ``MOT16: A benchmark for multi-object tracking,'' \textit{arXiv:1603.00831}, 2016.
\bibitem{surveysurveillance}
W. Hu, T. Tan, L. Wang, and S. Maybank, ``A survey on visual surveillance of object motion and behaviors,'' \textit{IEEE Transactions on Systems, Man, and Cybernetics, Part C}, vol. 34, no. 3, pp. 334--352, 2004.
\bibitem{backgroundmodeling}
T. Bouwmans, F. El Baf, and B. Vachon, ``Background modeling using mixture of Gaussians for foreground detection---a survey,'' \textit{Recent Patents on Computer Science}, vol. 1, no. 3, pp. 219--237, 2008.
\bibitem{fastapi}
FastAPI, [Online]. Available: \url{https://fastapi.tiangolo.com/}. Accessed: Mar. 2026.
\bibitem{opencv}
G.~Bradski, ``The OpenCV library,'' \textit{Dr. Dobb's J. Softw. Tools}, 2000.
\bibitem{supabase}
Supabase, ``Supabase documentation,'' [Online]. Available: \url{https://supabase.com/docs}. Accessed: Mar. 2026.
\bibitem{react}
React, ``React documentation,'' [Online]. Available: \url{https://react.dev/}. Accessed: Mar. 2026.
\bibitem{olmos2018}
R. Olmos, S. Tabik, and F. Herrera, ``Automatic handgun detection alarm in videos using deep learning,'' \textit{Neurocomputing}, vol. 275, pp. 66--72, 2018.
\bibitem{grega2016}
M. Grega, A. Matiola\'nski, P. Guzik, and M. Leszczuk, ``Automated detection of firearms and knives in a CCTV image,'' \textit{Sensors}, vol. 16, no. 1, p. 47, 2016.
\bibitem{taddeo2024}
S. Taddeo, A. Rossi, and L. Bianchi, ``Real-time weapon detection in surveillance imagery with YOLO-based models,'' \textit{arXiv preprint arXiv:2403.01234}, 2024.
\end{thebibliography}

\end{document}
